% !TEX root = ../../../main.tex
\section{Determinants of endomorphisms}
\label{sec:linear-algebra-i:determinants-of-endomorphisms}

% Sources: IMG_9263--IMG_9265, IMG_9317--IMG_9318;
% October 8 PDF, pp. 15--17.
An endomorphism is a linear map from a vector space to itself. If \(U\)
is one dimensional, every linear map \(F:U\to U\) is multiplication
by a scalar. For \(L:V\to V\), with \(\dim V=n\), apply this fact to
\[
  L^*:\Lambda^n V^*\longrightarrow\Lambda^n V^*.
\]

\begin{definition}
  The determinant of \(L:V\to V\) is the unique scalar \(\det L\)
  for which
  \begin{equation}
    L^*A=(\det L)A\qquad\text{for every \(A\in\Lambda^n V^*\)}.
    \label{eq:linear-algebra-i:determinant-pullback}
  \end{equation}
\end{definition}

Equivalently, for all \(v_1,\ldots,v_n\in V\),
\[
  A(L(v_1),\ldots,L(v_n))=(\det L)A(v_1,\ldots,v_n).
\]
The real two-dimensional case is visualized in
\autoref{fig:linear-algebra-i:determinant-linear-map}.
\begin{figure}[htbp]
  \centering
  % A real two-dimensional schematic of the already defined volume ratio.
\begin{tikzpicture}[
  x=11mm,y=11mm,font=\footnotesize,line width=0.45pt,
  vector/.style={-{Latex[length=1.6mm,width=1.1mm]}}
]
  \coordinate (O) at (0,0);
  \coordinate (V) at (2,0);
  \coordinate (W) at (0.4,1.3);
  \coordinate (S) at ($(V)+(W)$);
  \fill[black!8] (O) -- (V) -- (S) -- (W) -- cycle;
  \draw (V) -- (S) -- (W);
  \draw[vector] (O) -- (V) node[midway,below=3pt] {\(v\)};
  \draw[vector] (O) -- (W) node[midway,left=3pt] {\(w\)};
  \node at (1.1,-0.9) {\(A(v,w)\)};
  \draw[vector] (3.1,1) -- (4.4,1) node[midway,above=3pt] {\(L\)};
  \begin{scope}[xshift=64mm]
    \coordinate (O) at (0,0);
    \coordinate (V) at (2.1,0.4);
    \coordinate (W) at (0.6,1.8);
    \coordinate (S) at ($(V)+(W)$);
    \fill[black!15] (O) -- (V) -- (S) -- (W) -- cycle;
    \draw (V) -- (S) -- (W);
    \draw[vector] (O) -- (V) node[midway,below=3pt] {\(L(v)\)};
    \draw[vector] (O) -- (W) node[midway,left=3pt] {\(L(w)\)};
    \node at (1.3,-0.9) {\(A(L(v),L(w))\)};
  \end{scope}
\end{tikzpicture}

  \caption[Determinant as oriented-area scaling]{A real-plane schematic.
    Applying \(L\) to both edges multiplies their oriented area by
    \(\det L\). The drawing shows positive areas; a negative determinant
    reverses orientation.}
  \label{fig:linear-algebra-i:determinant-linear-map}
\end{figure}

Choose an ordered basis \(E=(e_1,\ldots,e_n)\) and normalize \(A\) so that
\[
  A(e_1,\ldots,e_n)=1.
\]
Then
\[
  \det L=A(L(e_1),\ldots,L(e_n)).
\]
For an unnormalized nonzero \(A\), the formula is
\[
  \det L=\frac{A(L(e_1),\ldots,L(e_n))}{A(e_1,\ldots,e_n)}.
\]

\subsection{Agreement with the representing matrix}

\begin{proposition}
  \label{thm:linear-algebra-i:determinant-representing-matrix}
  If \(L(E)=EM\), then \(\det L=\det M\).
\end{proposition}

\begin{proof}
  Write \(L(e_j)=\sum_{i=1}^n e_iM^i_j\), and normalize \(A\) on \(E\).
  Expanding each input gives
  \begin{align*}
    \det L
      &=A(L(e_1),\ldots,L(e_n))\\
      &=\sum_{k_1=1}^n\cdots\sum_{k_n=1}^n
        M^{k_1}_1\cdots M^{k_n}_n A(e_{k_1},\ldots,e_{k_n})\\
      &=\sum_{\sigma\in S_n}\varepsilon(\sigma)
        M^{\sigma(1)}_1\cdots M^{\sigma(n)}_n\\
      &=\det M.
  \end{align*}
  Repeated indices vanish, and each remaining basis value is the sign
  of its permutation, since \(A(e_1,\ldots,e_n)=1\).
\end{proof}

\subsection{Computations using multilinearity}

\begin{example}
  \label{ex:linear-algebra-i:triangular-endomorphism-determinant}
  Let \((e_1,e_2,e_3)\) be a basis and suppose
  \[
    \begin{aligned}
      L(e_1)&=2e_1+3e_2+4e_3,\\
      L(e_2)&=5e_2-6e_3,\\
      L(e_3)&=7e_3.
    \end{aligned}
  \]
  Choose \(A\) with \(A(e_1,e_2,e_3)=1\). Starting with the last
  input eliminates repeated basis vectors from the earlier inputs,
  \begin{align*}
    \det L
      &=A(2e_1+3e_2+4e_3,5e_2-6e_3,7e_3)\\
      &=A(2e_1+3e_2,5e_2,7e_3)\\
      &=A(2e_1,5e_2,7e_3)\\
      &=2\cdot5\cdot7\,A(e_1,e_2,e_3)=70.
  \end{align*}
  For example, the discarded term from the second input is
  \[
    -6A(L(e_1),e_3,7e_3)=0.
  \]
  After it is removed, the \(3e_2\)
  and \(4e_3\) terms in the first input also give repeated inputs.
\end{example}

% IMG_9265 has +2 at (4,2) and obtains -60. PDF p.17 has -2 and
% obtains +60. These are different matrices, not conflicting answers.
\begin{example}
  \label{ex:linear-algebra-i:sparse-determinant}
  Consider
  \[
    M=\begin{bmatrix}
      2&0&-1&0\\
      1&3&1&5\\
      0&0&3&0\\
      2&2&1&0
    \end{bmatrix}.
  \]
  Its columns are
  \[
    \begin{aligned}
      C_1&=2e_1+e_2+2e_4,& C_2&=3e_2+2e_4,\\
      C_3&=-e_1+e_2+3e_3+e_4,& C_4&=5e_2.
    \end{aligned}
  \]
  Let \(D\) be normalized on the standard basis. The last input
  eliminates the \(e_2\) terms, and the second input then eliminates
  the \(e_4\) terms. Therefore,
  \begin{align*}
    \det M
      &=D(C_1,C_2,C_3,C_4)\\
      &=D(2e_1+2e_4,2e_4,-e_1+3e_3+e_4,5e_2)\\
      &=D(2e_1,2e_4,-e_1+3e_3,5e_2)\\
      &=D(2e_1,2e_4,3e_3,5e_2)\\
      &=60D(e_1,e_4,e_3,e_2)\\
      &=-60D(e_1,e_2,e_3,e_4)=-60.
  \end{align*}
  The final sign comes from interchanging the second and fourth inputs.

  If the entry in row \(4\), column \(2\) is instead \(-2\), the
  same calculation gives
  \[
    -60D(e_1,e_4,e_3,e_2)=60.
  \]
\end{example}
